% `tex.web` §977 and §979: a `\vsplit` records the marks it took away.
%
%     if type(p)=mark_node then
%       if split_first_mark=null then
%         begin split_first_mark:=mark_ptr(p); split_bot_mark:=split_first_mark;
%         …
%       else begin delete_token_ref(split_bot_mark);
%         split_bot_mark:=mark_ptr(p); …
%
% The first `\mark` in the part taken off becomes `\splitfirstmark` and the
% last becomes `\splitbotmark`, and §977 clears both at the start of every
% `\vsplit` whether or not it finds any — so a split that takes no mark leaves
% them empty rather than leaving the last split's behind.
%
% We recorded neither. A document that splits a box and asks what was in it —
% which is what a footnote or a long table does to carry a heading over — got
% nothing back. In trip.log it is `3:alpha:3:alpha:alpha` against our
% `3:alpha:3::`.
\catcode`\{=1 \catcode`\}=2
\splitmaxdepth=2pt
\setbox1\vbox{\mark{A}\hbox{}\mark{B}\vskip 5pt\hbox{}\mark{C}\vskip 100pt\hbox{}\mark{D}}
\message{<before=[\splitfirstmark][\splitbotmark]>}
\setbox2\vsplit1 to 5pt
\message{<after=[\splitfirstmark][\splitbotmark]>}
\setbox3\vbox{\hbox{}}
\setbox4\vsplit3 to 5pt
% A split that finds no mark clears them rather than leaving the last one.
\message{<none=[\splitfirstmark][\splitbotmark]>}
\end
